\documentclass[11pt]{article}
\usepackage[margin=1in]{geometry}
\usepackage{amsmath,amssymb,amsthm}
\usepackage{graphicx}
\usepackage{microtype}
\usepackage{xcolor}
\usepackage{enumitem}
\usepackage[colorlinks=true,linkcolor=blue!55!black,citecolor=blue!55!black,urlcolor=blue!55!black]{hyperref}

\newtheorem{theorem}{Theorem}
\newtheorem{corollary}[theorem]{Corollary}
\newcommand{\R}{\mathbb R}
\newcommand{\Z}{\mathbb Z}
\newcommand{\N}{\mathbb N}

\title{Lattice balls refute the exponent restriction in Conjecture 4.2\\
\large A counterexample to arXiv:1412.6400 as printed}
\author{Automated Mathematics Run \texttt{fa\_banach\_001}, lane 12}
\date{9 August 2026}

\begin{document}
\maketitle

\begin{center}
\fcolorbox{blue!55!black}{blue!4}{\parbox{0.89\textwidth}{
\textbf{Status.} Candidate full counterexample, subject to human
verification.  In dimension $3$, the polynomial
\[
 P(x)=1+x_1^2+x_2^2+x_3^2
\]
has finite lattice sublevel sets with cardinality comparable to $t^{3/2}$.
Conjecture 4.2 permits only powers $t^{\alpha/\beta}$ with
$0<\alpha\le\beta$, hence exponent at most $1$.}}
\end{center}

\section{Source and exact conjecture}

The source is P.~L.~Combettes and D.~D\~ung,
\emph{Kolmogorov $n$-Widths of Function Classes Induced by a Non-Degenerate
Differential Operator: A Convex Duality Approach}~\cite{CD},
arXiv:1412.6400, published in \emph{Set-Valued and Variational Analysis}
24 (2016), 83--99.

For a real polynomial $P$ on $\R^d$, the source defines
\[
 K(t)=\{k\in\Z^d:|P(k)|\le t\}.
\]
Conjecture 4.2, in Section 4.6 on source PDF page 19, states:

\begin{quote}
If $K(t)$ is finite for every $t\ge0$, then there exist integers
$\alpha,\beta,\nu$ with
\[
 0<\alpha\le\beta,\qquad 0\le\nu<d,
\]
such that, for all sufficiently large $t$,
\begin{equation}\label{eq:conjecture}
 \#K(t)\asymp t^{\alpha/\beta}(\log t)^\nu.
\end{equation}
\end{quote}
Here $A(t)\asymp B(t)$ means that their ratio stays between two positive
constants for all sufficiently large $t$.  Figure~\ref{fig:conjecture}
shows the exact printed statement.

\begin{figure}[ht]
 \centering
 \includegraphics[width=\textwidth]{figures/open_problem_crop.png}
 \caption{Conjecture 4.2 on page 19 of the source paper~\cite{CD}.}
 \label{fig:conjecture}
\end{figure}

The restriction $\alpha\le\beta$ forces $\alpha/\beta\le1$.  This is
violated already by ordinary lattice balls.

\section{Proof intuition}

Use a positive coercive polynomial whose sublevel sets are dilates of a
fixed full-dimensional body.  In $d$ dimensions, a degree-$m$ polynomial
of the form
\[
 1+x_1^m+\cdots+x_d^m\qquad(m\ \text{even})
\]
has sublevel sets of radius $t^{1/m}$.  A full-dimensional lattice ball of
that radius contains and is contained in coordinate cubes of comparable
side length, so it has about $(t^{1/m})^d=t^{d/m}$ lattice points.  Taking
$d>m$ makes the power larger than one, while every sublevel set remains
finite.  Thus the example satisfies the sole hypothesis of the conjecture
and violates its permitted conclusion.

\section{Counterexample}

\begin{theorem}[An infinite counterexample family]\label{thm:family}
Let $m\ge2$ be even and let $d>m$.  Define
\[
 P_{d,m}(x_1,\ldots,x_d)=1+\sum_{j=1}^d x_j^m.
\]
Then $K_{d,m}(t)=\{k\in\Z^d:|P_{d,m}(k)|\le t\}$ is finite for every
$t\ge0$, but
\begin{equation}\label{eq:growth}
 \#K_{d,m}(t)\asymp t^{d/m}.
\end{equation}
Consequently $P_{d,m}$ refutes Conjecture 4.2 as printed.
\end{theorem}

\begin{proof}
Because $m$ is even,
\[
 P_{d,m}(k)=1+\sum_{j=1}^d |k_j|^m\ge1.
\]
Thus $K_{d,m}(t)$ is empty for $t<1$, while for $t\ge1$ it is the
lattice $\ell_m$-ball
\begin{equation}\label{eq:ball}
 K_{d,m}(t)=\left\{k\in\Z^d:
       \sum_{j=1}^d|k_j|^m\le t-1\right\}.
\end{equation}
In particular every such set is finite.

Put $R=(t-1)^{1/m}$.  The elementary inclusions
\begin{equation}\label{eq:cubes}
 \left[-\frac{R}{d^{1/m}},\frac{R}{d^{1/m}}\right]^d
 \subseteq\left\{x\in\R^d:\sum_{j=1}^d|x_j|^m\le R^m\right\}
 \subseteq[-R,R]^d
\end{equation}
give, after counting integer points,
\[
 \left(2\left\lfloor\frac{R}{d^{1/m}}\right\rfloor+1\right)^d
 \le \#K_{d,m}(t)
 \le (2\lfloor R\rfloor+1)^d.
\]
For sufficiently large $t$, both outer expressions are constant multiples
of $R^d$.  Hence
\[
 \#K_{d,m}(t)\asymp R^d=(t-1)^{d/m}\asymp t^{d/m},
\]
proving \eqref{eq:growth}.

Suppose the conclusion of Conjecture 4.2 held.  Set
$r=\alpha/\beta$.  Since $0<\alpha\le\beta$, we have $0<r\le1$;
also $0\le\nu<d$.  But $d/m>1$, and therefore
\[
 \frac{\#K_{d,m}(t)}{t^r(\log t)^\nu}
 \asymp\frac{t^{d/m-r}}{(\log t)^\nu}\longrightarrow\infty.
\]
This contradicts the upper comparison required by
\eqref{eq:conjecture}.
\end{proof}

\begin{corollary}[Smallest simple instance]
For
\[
 d=3,\qquad P(x_1,x_2,x_3)=1+x_1^2+x_2^2+x_3^2,
\]
every $K(t)$ is finite and $\#K(t)\asymp t^{3/2}$.  Hence this single
quadratic polynomial refutes Conjecture 4.2.
\end{corollary}

\section{Verification, interpretation, and novelty}

No external theorem or numerical computation is used.  The hypothesis is
verified directly by coercivity, and the two cube inclusions
\eqref{eq:cubes} prove the exact growth order.  A reviewer only needs to
check that Conjecture 4.2 indeed includes the printed restriction
$0<\alpha\le\beta$; both the arXiv PDF and the authors' currently available
published-version PDF contain that restriction.

The counterexample is deliberately scoped to the statement as printed.  If
the intended conjecture merely asked for an arbitrary positive rational
power, with no requirement $\alpha\le\beta$, this family would satisfy the
corrected conclusion by taking $\alpha=d$, $\beta=m$, and $\nu=0$.
Therefore the packet does not refute that natural corrected conjecture.

A bounded search through 9 August 2026 used arXiv:1412.6400, the exact
paper title, the authors' names, ``Conjecture 4.2'', ``card $K(t)$'', and
the displayed $\alpha/\beta$ formula.  It located the arXiv record, the
2016 publication metadata, and the authors' published-version PDF, which
still prints the same inequality.  No erratum, later correction, or paper
stating this counterexample was found.  Novelty is plausible, but the
most likely interpretation is a persistent typographical restriction rather
than a failure of the broader lattice-counting program.

\paragraph{Human-review recommendation.}
Verify whether the authors intended to omit $\alpha\le\beta$ or to reverse
it.  As a mathematical statement, Conjecture 4.2 in both located versions
is false by Theorem~\ref{thm:family}.

\begin{thebibliography}{9}
\bibitem{CD}
P.~L.~Combettes and D.~D\~ung,
\emph{Kolmogorov $n$-Widths of Function Classes Induced by a Non-Degenerate
Differential Operator: A Convex Duality Approach},
Set-Valued and Variational Analysis \textbf{24} (2016), 83--99;
arXiv:1412.6400.
\end{thebibliography}

\end{document}
